\documentclass[sigconf,nonacm]{acmart}

\usepackage{booktabs}

\settopmatter{printacmref=false}
\renewcommand\footnotetextcopyrightpermission[1]{}
\pagestyle{plain}

\begin{document}

\title{Composable Exact Bounds: A Self-Configuring Hybrid of PCA and
Triangle-Inequality Pruning}

\author{Tai-Wu Chiang}
\email{twchiang@utexas.edu}
\affiliation{%
  \institution{Independent Researcher}
  \city{California}
  \country{USA}}

\begin{abstract}
Exact nearest-neighbor search can be accelerated by cheap lower bounds that
prune candidates without computing full distances. Our prior work used bounds
from truncated PCA projections, which are tight precisely when variance
concentrates in few dimensions. We observe that such bounds are one member of
a \emph{family}: any region that provably contains a set of points yields an
exact lower bound via the query-to-region distance, and different region
shapes are tight for different data structures. We compose two complementary
families---PCA truncation (tight for linearly correlated data) and
triangle-inequality ball bounds around $k$-means centroids (tight for
clustered data, and exact for \emph{any} centers, so approximation quality
affects only speed, never correctness)---into a single exact cascade. Because
neither family helps on all data, the index \emph{configures itself}: a
three-tier diagnostic measures a sample's eigenvalue decay, detects structure
beyond the covariance with two cheap statistics, and settles the question
with a measured trial of the real cascade. We report results in
hardware-independent work units---full-vector fetches, total bytes touched
(charging the pruning machinery itself), and resident RAM per
vector---because wall-clock proved regime-dependent: the same algorithm
measured $4.7\times$, $3.6\times$, $2.7\times$, and parity under changes of
numeric type, machine load, and implementation tuning, while its work counts
never moved. On SIFT1M with official queries, the cascade computes full distances
for 0.06\% of the dataset and keeps 8--40 bytes per vector resident in RAM
(versus 512 for a flat scan); on clustered data with a flat eigenvalue
spectrum, the ball level prunes to 3\% before a single projection is computed
($13.9\times$ wall-clock versus $2.1\times$ for PCA alone); on SIFT the ball
level's bound work cancels its savings---a verdict the diagnostic correctly
predicts from a sample, protecting the practitioner from building a useless
level. Exactness is verified against ground truth on every query. All code,
diagnostics, and experiments are public.
\end{abstract}

\maketitle

\section{Introduction}

Finding the exact nearest neighbor in a large, high-dimensional dataset is
bottlenecked by the cost of computing full distances. Approximate methods
(LSH, HNSW~\cite{hnsw}, IVF~\cite{faiss}) trade recall for speed; for
applications where a missed true neighbor is unacceptable, the alternative is
exact pruning: compute a cheap \emph{lower bound} on each candidate's
distance, and discard the candidate if the bound already exceeds the best
distance found so far. Our prior work~\cite{paper1} built such an index from
truncated PCA projections: the distance between $k$-dimensional PCA prefixes
lower-bounds the true distance, so a cascade of progressively finer prefixes
prunes over 99.9\% of the SIFT1M dataset while returning provably exact
results.

That bound has a blind spot, and it is structural. PCA truncation is tight
exactly when variance concentrates in few principal directions---linearly
correlated data. Datasets whose structure is \emph{clustered} rather than
correlated can present a flat eigenvalue spectrum, collapsing the PCA
cascade toward brute force, while remaining highly compressible in a
different sense: points huddle near a modest number of centroids. No single
bound family serves all data, and, worse, a practitioner has no way to know
before building an index which family---if any---their data supports.

This paper makes two moves. First, \textbf{exact bounds compose}: the
maximum of valid lower bounds is a valid lower bound, so bound families with
complementary blind spots can share one cascade. We add a
triangle-inequality \emph{ball bound} around $k$-means centroids---the
clustered-data complement of PCA truncation---and show that it is exact for
\emph{any} choice of centers: clustering quality affects only pruning power,
never correctness. Second, \textbf{the composition is chosen by
measurement}: a three-tier diagnostic, run on a sample in seconds, rates the
PCA levels from the eigenvalue decay, detects structure beyond the
covariance with two cheap statistics, and---because we show statistics alone
mispredict in both directions---settles the question with a measured trial
of the real cascade, emitting an \textsc{add}/\textsc{skip} verdict per
level.

There is also an economic motivation. An exhaustive flat scan requires the
entire dataset resident in RAM---the most expensive byte in the machine,
running roughly an order of magnitude above NVMe per byte, and more under
supply pressure. The pruning cascade needs only its metadata resident: 8--40
bytes per vector versus 512 bytes per vector for \texttt{fp32} data. At one
billion vectors this is 8--40\,GB of RAM versus 512\,GB---a commodity box
versus a specialty one---with exactness preserved.

Our contributions:
\begin{enumerate}
\item A \textbf{containing-region principle} unifying exact pruning bounds,
  with the ball bound as the clustered-data complement to PCA truncation
  (\S\ref{sec:theory}).
\item A \textbf{composed exact cascade}---ball level plus PCA levels---with a
  radius re-probe that keeps the cascade as tight as PCA-only. Storage
  overhead is $\sim$8 bytes per point (\S\ref{sec:algorithm}).
\item A \textbf{three-tier self-configuration funnel}: the statistics
  nominate, the measurement decides; negative verdicts included
  (\S\ref{sec:funnel}).
\item \textbf{Reproducible evidence in hardware-independent work units},
  with wall-clock demoted to validating a cost model, plus negative
  controls and a sample-to-full-scale extrapolation check
  (\S\ref{sec:eval}).
\end{enumerate}

\section{Background and Related Work}

\textbf{Paper I in brief.} For data rotated onto its principal axes, the
squared distance restricted to the first $k$ coordinates is a lower bound on
the full squared distance (the discarded terms are non-negative). The
algorithm of~\cite{paper1} exploits this with a cascade of prefix widths
$(8, 16, 32, 64)$ and a dynamic radius: a Top-50 probe at the coarsest level
computes exact distances for the 50 most promising candidates, seeding a
mathematically safe pruning radius without prior knowledge. On SIFT1M the
cascade discards all but an average of 564 of 1{,}000{,}000 vectors before
any full-dimensional evaluation.

\textbf{Lineage.} Neither bound family is new, and we do not claim
otherwise. Lower-bounding by dimensionality reduction descends from the
GEMINI framework~\cite{gemini} and the VA-file~\cite{vafile}. Pruning by
the triangle inequality around reference points is the basis of ball
trees~\cite{balltree}, vantage-point trees~\cite{vptree}, GNAT~\cite{gnat},
and Elkan's accelerated $k$-means~\cite{elkan}. The approximate cousins are
FAISS IVF~\cite{faiss} (the same $k$-means geometry without exactness) and
OPQ~\cite{opq} (a per-cell rotation, the ellipsoid rung of the shape ladder
in \S\ref{sec:theory}). Disk-resident ANN systems such as
DiskANN~\cite{diskann} and SPANN~\cite{spann} target the same
memory-economics regime we do, but are approximate. Our contribution is the
\emph{exact composition} of the two bound families in one cascade, the
\emph{measured self-configuration}, and a diagnostic that also tells the
practitioner when to walk away.

\section{Composable Exact Bounds}
\label{sec:theory}

\subsection{The Containing-Region Principle}

Let $R \subseteq \mathbb{R}^d$ be any region that provably contains a set of
points $C$. Then for every $p \in C$ and any query $q$,
\begin{equation}
\lVert q - p \rVert \;\ge\; \mathrm{dist}(q, R),
\end{equation}
so $\mathrm{dist}(q,R)$ is an exact pruning bound for every member of $C$
at once. A bound family is therefore a \emph{shape} family, and a shape is
admissible when (a) containment is guaranteed and (b) $\mathrm{dist}(q,R)$
is cheap. Table~\ref{tab:shapes} lists the natural ladder. Each geometric
family is the twin of a quantizer family: a quantizer's
codeword-plus-distortion model \emph{is} a containing region. We use the
first two rungs; the others are available under the same principle.

\begin{table}[t]
\caption{Shape families and their exact bounds. Each is tight for a
different data structure; each mirrors a quantizer family.}
\label{tab:shapes}
\small
\begin{tabular}{@{}llll@{}}
\toprule
Shape & Bound cost & Tight for & Quantizer twin \\
\midrule
ball & $|\,\lVert q{-}c\rVert - d_p\,|$ & clusters & VQ \\
ellipsoid & per-axis scaled & stretched clusters & OPQ \\
slab & projection residual & local low rank & local PCA \\
cone$\times$annulus & closed form & radial/gain data & gain--shape \\
\bottomrule
\end{tabular}
\end{table}

\subsection{The Ball Bound}

Cluster the dataset with $k$-means. For each point $p$ assigned to centroid
$c$, store the scalar $d_p = \lVert p - c \rVert$ once, offline. At query
time, compute the $k$ centroid distances $\lVert q - c_j \rVert$ (with
$k \ll N$), and the triangle inequality gives, for every point,
\begin{equation}
\label{eq:ball}
\lVert q - p \rVert \;\ge\; \bigl|\, \lVert q - c_{a(p)} \rVert - d_p \,\bigr|,
\end{equation}
where $a(p)$ is $p$'s cluster. Evaluating (\ref{eq:ball}) for all $N$ points
requires one gather and one subtraction per point---\emph{no distance
computations}. A coarser, per-cluster form uses the cluster radius
$R_j = \max_{p \in j} d_p$: if $\lVert q - c_j \rVert - R_j$ exceeds the
current radius, the entire cluster is eliminated in one comparison, which in
an out-of-core layout skips a whole contiguous block (\S\ref{sec:layout}).

\subsection{Exactness Is Unconditional; Quality Is Not}

Inequality~(\ref{eq:ball}) is the triangle inequality: it holds for
\emph{any} reference points whatsoever. $k$-means enters only because it
minimizes the mean $d_p^2$, which is what makes the bound \emph{tight}. Three
corollaries follow. Stale centers remain correct as the dataset drifts
(pruning degrades gracefully; re-cluster when a diagnostic says so).
Insertion is $O(1)$: assign the new point to its nearest existing centroid
and store $d_p$; no re-clustering is required for correctness. And the
composed index inherits Paper~I's property in mirror image: there, any
orthogonal rotation is exact and PCA makes it tight; here, any centers are
exact and $k$-means makes them tight. Since the maximum of lower bounds is a
lower bound, the two families compose freely in one cascade.

\subsection{Complementarity}

PCA truncation is tight iff the eigenvalue spectrum decays steeply---a
\emph{global, linear} property. The ball bound is tight iff points sit close
to centroids relative to the spread of $\lVert q - c \rVert$---a
\emph{local} property that holds for multimodal data and, more generally,
for data of low intrinsic dimension, where $d_p$ and the centroid distances
spread widely instead of concentrating. The blind spots are disjoint by
construction. Our clearest demonstration: a synthetic dataset of 512 tight
clusters whose centers span all dimensions has a nearly flat spectrum (95\%
of variance needs 111 of 128 dimensions; the PCA cascade manages
$2.1\times$), yet the ball bound prunes to 3\% before a single projection is
computed.

\section{The Self-Configuration Funnel}
\label{sec:funnel}

The funnel answers, from a sample and in seconds, ``which levels should this
index have for \emph{my} data?''

\textbf{Tier 1 --- eigenvalue decay.} Fit PCA on a sample; the cumulative
explained variance at each cascade width is the expected tightness of that
PCA level. This tier rates the PCA cascade (and reproduces Paper~I's SIFT1M
behavior from a 100k sample).

\textbf{Tier 2 --- beyond-covariance statistics.} Two cheap statistics on
the PCA-rotated sample: the mean excess kurtosis of components (marginal
non-Gaussianity), and $\mathrm{sqdep}$, the mean
$|\mathrm{corr}(z_i^2, z_j^2)|$ over component pairs (joint dependence that
no rotation can remove; finite-sample noise floor $\approx\sqrt{2/n}$).
These detect that structure beyond the covariance \emph{exists}---not
whether it is \emph{harvestable}.

\textbf{Tier 3 --- the measured trial.} Fit the real hybrid index on the
sample, run the real cascade with the ball level on and off, and report the
measured ratio and the ball stage's own pruning rate. The verdict is
\textsc{add} or \textsc{skip}.

Tier 3 must be the decider, and the evidence is that the statistics
mispredict in \emph{both} directions. On SIFT, tier 2 fires
($\mathrm{sqdep} = 0.032$, $7\times$ the floor: real beyond-covariance
structure) yet the trial says \textsc{skip} ($1.0\times$): ball bounds prune
only 50--59\%, and the bound work cancels the savings. On a
latent-dimension-8 Gaussian, tier 2 is clean (kurtosis $\approx 0$,
$\mathrm{sqdep}$ at floor) yet the trial says \textsc{add} ($1.9\times$):
low intrinsic dimension spreads $d_p$ widely, so the bounds bite (pruning to
23\%) with zero multimodality. Across all datasets the measured value of the
ball level is a monotone function of its own pruning rate---survivors of
3\% $\rightarrow$ $4$--$14\times$; 15--23\% $\rightarrow$ $2$--$3\times$;
50--59\% $\rightarrow$ parity; 100\% $\rightarrow$ $0.8\times$ (a measured
loss)---and the trial measures exactly this quantity. No statistic
substitutes for it.

\section{The Hybrid Algorithm}
\label{sec:algorithm}

\subsection{Offline}

Fit PCA (as in Paper~I) and $k$-means with
$k = \mathrm{clip}(\sqrt{N}, 64, 4096)$ (MiniBatch; 24\,s on SIFT1M for
$k{=}1024$). Store per point a cluster id (\texttt{int32}) and $d_p$
(\texttt{float32})---8 bytes---plus the $k \times d$ centroids and
per-cluster radii.

\subsection{Query Pipeline}

\begin{enumerate}
\item $k$ centroid distances ($k \cdot d$ operations).
\item Ball bounds (\ref{eq:ball}) for all $N$ points: one gather and one
  absolute difference each.
\item \textbf{Probe}: exact distances to the \textsc{probe}$=50$ points with
  the smallest ball bounds seed a safe initial radius $r$ (Paper~I's Top-50
  probe, seeded by ball bounds instead of coarse projections).
\item Ball pruning: retain points with $\mathrm{bound}^2 < r^2$.
\item PCA cascade $(8,16,32,64)$ on the survivors, with a \textbf{re-probe}
  at the first PCA stage: before pruning, exact distances to the 50 best
  8-dimensional candidates tighten $r$. Without this step the ball-seeded
  radius is loose and the cascade degrades (8-D survivors 27.2\% versus
  14.6\% with it); with it, the cascade matches PCA-only tightness
  stage-for-stage.
\item Full-distance check on the final survivors. The result is provably
  identical to brute force, and is additionally verified against the
  official ground truth in every experiment.
\end{enumerate}

\subsection{Out-of-Core Layout}
\label{sec:layout}

Pruning fragments I/O only if the layout ignores it---the classic
index-versus-scan crossover: 600 scattered 512-byte reads lose to a full
sequential scan on seek-bound media. Measured on SIFT1M ($k{=}1024$,
$\sim$488\,KB blocks): the cluster-level bound alone is a weak block-skipper
(40.4\% of clusters survive with a 95th-percentile radius and a 5\% escape
set; 52.3\% with the max radius), \emph{but} the final $\sim$600 survivors
concentrate in a median of 44 of 1024 clusters. A cluster-major layout
therefore converts the fetch into $\sim$44 contiguous reads
($\sim$21\,MB). The resulting two-tier design: (1) scan the small,
sequential metadata---\texttt{assign}${+}d_p$ is 8\,MB, and optionally the
8-dimensional PCA sketch, 32\,MB, i.e.\ $64\times$/$16\times$ smaller than
the data; pruning decisions never touch full vectors---then (2) fetch
survivors as batched random 512-byte reads on NVMe ($\sim$0.3\,MB) or as the
$\sim$44 cluster blocks on seek-bound media. Back-of-envelope: NVMe
$\sim$15\,ms versus 146\,ms for a full scan ($\sim$10$\times$); HDD
$\sim$0.55\,s versus 3.4\,s ($\sim$6$\times$), where the naive scattered
layout would \emph{lose} ($\sim$4.8\,s). This is the wisdom of the
VA-file~\cite{vafile}, SPANN~\cite{spann}, and DiskANN~\cite{diskann}:
sequential sketches, layout-aware fetches. The same $k$-means that supplies
the bound supplies the layout.

\section{Evaluation}
\label{sec:eval}

\textbf{Metrics policy.} Primary claims are stated in hardware-independent
work units, because wall-clock proved regime-dependent: the same algorithm
measured $4.7\times$, $3.6\times$, $2.7\times$, and parity under changes of
numeric type, machine load, and implementation tuning, while its work counts
never moved (details below). The primary metrics are (1) full-vector fetches
per query (the out-of-core I/O currency), (2) total bytes touched per query,
\emph{charging the pruning machinery itself}, (3) resident RAM bytes per
vector, and (4) the survivor cascade, which is the algorithmic fingerprint.
Wall-clock appears as validation of a first-order cost model
($t \approx \mathrm{bytes}/\mathrm{bandwidth} +
\mathrm{ops}/\mathrm{throughput}$), with FAISS supplying the tuned-constants
reference---so a reader can plug their machine's constants into our work
numbers and predict their own outcome.

\textbf{Setup.} Intel i9-9900K (8C/16T, AVX2${+}$FMA, no AVX-512),
dual-channel DDR4; \texttt{faiss-cpu} 1.15 dispatching its AVX2 build; NumPy
and OpenBLAS at the AVX2 tier---baseline and FAISS use the same instruction
set. Timings are medians of 5 repetitions on a quiet machine
with warm cache, exactness verified against the official ground truth on
every repetition. Results are CPU-governor-insensitive (balanced versus
performance profiles: all medians within 2\%). Spreads are $\le$2\% except
where noted.

\subsection{SIFT1M}

\begin{table}[t]
\caption{SIFT1M, 100 official queries, medians of 5 repetitions,
implementation parity. All methods return the ground-truth neighbor on
every query.}
\label{tab:sift}
\small
\begin{tabular}{@{}lrrl@{}}
\toprule
Method & ms/query & Speedup & Survivors per stage (\%) \\
\midrule
brute (NumPy) & 167.3 & 1.0$\times$ & --- \\
PCA-only~\cite{paper1} & 61.1 & 2.7$\times$ & 14.7 / 5.2 / 0.8 / 0.06 \\
hybrid & 62.2 & 2.7$\times$ & ball 50.6 $\to$ 14.6 / 5.1 / 0.7 / 0.05 \\
\midrule
FAISS-flat, 1 thread & 29.3 & 5.7$\times$ & (scans 100\%) \\
FAISS-flat, batch-100 & 5.9 & 28$\times$ & (scans 100\%) \\
\bottomrule
\end{tabular}
\end{table}

Table~\ref{tab:sift} gives the in-RAM picture. The PCA cascade alone reaches
$2.7\times$; \textbf{the ball level is a wall-clock tie on SIFT}---its
50\% pruning saves about what its bound computation costs---exactly as the
sample trial predicts (\textsc{skip}, \S\ref{sec:funnel}). The hybrid's
SIFT-side value is therefore out-of-core and residency only
(\S\ref{sec:layout}); the wall-clock case for the ball level rests on
clustered and low-intrinsic-dimension data (\S\ref{sec:regimes}). The
survivor cascade reproduces Paper~I exactly.

FAISS \texttt{IndexFlatL2} wins the in-RAM race outright, scanning all
$10^6$ vectors in 29.3\,ms single-threaded. The gap is fusion, not SIMD:
both sides run AVX2, but FAISS's kernel makes one pass ($\sim$512\,MB of
traffic per query) where NumPy's expression materializes temporaries
($\sim$3 passes, 1.5--2\,GB), a $5.6\times$ advantage at identical work. The
comparison scopes the claim rather than defeating it: the cascade touches
$1{,}700\times$ fewer full vectors, requires 8--40 resident bytes per vector
against 512 for any flat scan, and its bounds compose with FAISS-grade
kernels rather than competing with them (the cascade could drive such a
kernel over its survivor sets). Disk-resident ANN systems occupy the same
economics but surrender exactness.

\textbf{Why wall-clock is the wrong headline (measured).} Two confounds were
caught only by interrogating ratio changes. \emph{Numeric type}: sklearn's
PCA silently outputs \texttt{float64}, whose bandwidth-bound brute baseline
is $1.86\times$ costlier; a $2\times2$ decomposition on identical code
measures the same cascade at $3.6\times$ (f64---Paper~I's regime, matching
its reported $4.3\times$) versus $2.7\times$ (f32). \emph{Implementation
parity}: an unnecessary full-array gather in our own PCA-only path inflated
the ball level's apparent contribution to $2.0\times$; fixing it revealed
the tie in Table~\ref{tab:sift}. Machine load adds a further multiplier:
background bandwidth contention slows the scan-bound baseline most,
flattering all pruners. Survivor counts were invariant through every one of
these changes. Hence the metrics policy.

\subsection{Structure Regimes}
\label{sec:regimes}

\begin{table}[t]
\caption{Structure regimes (synthetic, $n=200$k, $d=64$ for the first two;
$n=60$k, $d=128$ for the funnel rows), implementation parity. ``Ball
survivors'' is the ball stage's own pruning rate---the quantity the funnel's
trial measures.}
\label{tab:regimes}
\small
\begin{tabular}{@{}lrrrl@{}}
\toprule
Dataset & PCA-only & Hybrid & Ball surv. & Verdict \\
\midrule
512 tight clusters & 2.1$\times$ & \textbf{13.9$\times$} & 3\% & \textsc{add} \\
latent-8 Gaussian & 2.1$\times$ & 5.6$\times$ & 15--23\% & \textsc{add} \\
SIFT1M & 2.7$\times$ & 2.7$\times$ & 50--59\% & \textsc{skip} \\
latent-48 Gaussian & 1.9$\times$ & 0.8$\times$ of PCA & 100\% & \textsc{skip} \\
i.i.d.\ Gaussian & 0.3$\times$ & --- & 100\% & \textsc{skip} (all) \\
\bottomrule
\end{tabular}
\end{table}

Table~\ref{tab:regimes} spans the design space, negative results included
deliberately. On the flat-spectrum clustered dataset---PCA's blind
spot---the ball level prunes to 3\% before any projection is computed and
the hybrid reaches $13.9\times$; its bound requires no multiplications, so
it undercuts even the 8-dimensional scan. On the latent-48 blob the ball
level is measured to \emph{hurt} ($0.8\times$: 100\% survivors, pure
overhead), and on i.i.d.\ noise the diagnostic correctly reports that
nothing works and brute force or ANN should be used. The i.i.d.\ row also
shows the PCA cascade itself losing to brute force ($0.3\times$)---the
honest boundary of the entire approach.

\subsection{Extrapolation}

The funnel's trial on a 100k SIFT sample returns $1.0\times$
(\textsc{skip}); the full-1M parity benchmark measures exactly that (61.1
versus 62.2\,ms). The diagnostic predicts full-scale outcomes from a
sample---including the negative verdict, which is its most valuable
behavior: it prevents building a useless level.

\section{Limitations}

High \emph{intrinsic} dimensionality defeats every cheap bound: distance
concentration flattens the eigenvalue spectrum and the spread of $d_p$
simultaneously, and the funnel exists precisely to detect this and say so
(Table~\ref{tab:regimes}, last row). In-RAM wall-clock at native-code
constants belongs to fused exhaustive scans (FAISS) on this dataset scale;
our wall-clock case is confined to data the ball level can prune hard, and
to out-of-core and RAM-constrained regimes, which we argue from measured
work counts and a cost model rather than end-to-end I/O runs (a cold-cache
wall-clock run is future work). The current implementation is 1-NN;
top-$k$ generalizes the radius ($r$ = current $k$-th best) but is
unimplemented. Ball-bound quality depends on the sample used for $k$-means
in the same way any learned structure does, though correctness never does.

\section{Conclusion}

Exact lower bounds are shapes; shapes compose; and the right composition is
a property of the data, measurable on a sample before any index is built.
The resulting artifact---a hybrid of PCA truncation and triangle-inequality
ball pruning that configures itself, states its results in
hardware-independent work units, and honestly reports where it does not
help---extends our prior exact-search work to clustered data, cuts resident
RAM by $12$--$64\times$, and is available, with every experiment in this
paper, at the repository below. Measure the data to choose the index; count
the work to predict the time.

\vspace{1ex}
\noindent\textbf{Artifact.} Code, diagnostics, and all experiment scripts:\\
\texttt{github.com/taiwuchiang-gmail/hierarchical-pca-search}

\begin{thebibliography}{13}

\bibitem{paper1}
T.-W. Chiang.
\newblock A hierarchical pruning algorithm for fast, exact nearest neighbor
  search in high-dimensional spaces.
\newblock Zenodo, 2026. doi:10.5281/zenodo.21387359.

\bibitem{gemini}
C. Faloutsos, M. Ranganathan, and Y. Manolopoulos.
\newblock Fast subsequence matching in time-series databases.
\newblock In {\em SIGMOD}, 1994.

\bibitem{vafile}
R. Weber, H.-J. Schek, and S. Blott.
\newblock A quantitative analysis and performance study for similarity-search
  methods in high-dimensional spaces.
\newblock In {\em VLDB}, 1998.

\bibitem{pq}
H. J{\'e}gou, M. Douze, and C. Schmid.
\newblock Product quantization for nearest neighbor search.
\newblock {\em IEEE TPAMI}, 33(1), 2011.

\bibitem{opq}
T. Ge, K. He, Q. Ke, and J. Sun.
\newblock Optimized product quantization.
\newblock In {\em CVPR}, 2013.

\bibitem{elkan}
C. Elkan.
\newblock Using the triangle inequality to accelerate $k$-means.
\newblock In {\em ICML}, 2003.

\bibitem{balltree}
S. M. Omohundro.
\newblock Five balltree construction algorithms.
\newblock Technical report, ICSI, 1989.

\bibitem{vptree}
P. N. Yianilos.
\newblock Data structures and algorithms for nearest neighbor search in
  general metric spaces.
\newblock In {\em SODA}, 1993.

\bibitem{gnat}
S. Brin.
\newblock Near neighbor search in large metric spaces.
\newblock In {\em VLDB}, 1995.

\bibitem{hnsw}
Y. A. Malkov and D. A. Yashunin.
\newblock Efficient and robust approximate nearest neighbor search using
  hierarchical navigable small world graphs.
\newblock {\em IEEE TPAMI}, 42(4), 2018.

\bibitem{faiss}
J. Johnson, M. Douze, and H. J{\'e}gou.
\newblock Billion-scale similarity search with GPUs.
\newblock {\em IEEE Transactions on Big Data}, 2019.

\bibitem{diskann}
S. Jayaram Subramanya et al.
\newblock DiskANN: Fast accurate billion-point nearest neighbor search on a
  single node.
\newblock In {\em NeurIPS}, 2019.

\bibitem{spann}
Q. Chen et al.
\newblock SPANN: Highly-efficient billion-scale approximate nearest neighbor
  search.
\newblock In {\em NeurIPS}, 2021.

\end{thebibliography}

\end{document}
